\documentclass[12pt,a4paper]{report}
\usepackage[utf8]{inputenc}
\usepackage[T1]{fontenc}
\usepackage[english]{babel}
\usepackage[a4paper, left=2.5cm, right=2.5cm, top=3cm, bottom=2.5cm]{geometry}
\usepackage{graphicx}
\usepackage{caption}
\usepackage{amsmath, amssymb}
\usepackage{float}
\usepackage{xcolor}
\usepackage[table]{xcolor}

\usepackage{tikz}
\usetikzlibrary{shapes.geometric, arrows.meta}
\usepackage{pdfpages}
\usepackage{fancyhdr}
\usepackage{tocloft}
\usepackage[hidelinks]{hyperref}
\usepackage{enumitem}
%\usepackage{sansmath}
%\sansmath
%\renewcommand{\familydefault}{\sfdefault}
\usepackage{parskip}
\setlength{\parskip}{12pt}
\renewcommand{\cftsecleader}{\cftdotfill{\cftdotsep}}
\renewcommand{\arraystretch}{1.2}  

%\setlist[itemize]{label=-}
\setlist[itemize]{label=-, itemsep=0pt, parsep=0pt}



% COVER background
\usepackage{eso-pic}
\newcommand\BackgroundPic{
    \put(0,0){
        \parbox[b][\paperheight]{\paperwidth}{
            \vfill
            \centering
            \includegraphics[width=\paperwidth,height=\paperheight]{image/cover.pdf}
            \vfill
        }
    }
}

\begin{document}

%===== add cover ======
\AddToShipoutPicture*{\BackgroundPic} 
\thispagestyle{empty} 
% Chèn text lên nền
\vspace*{2cm}

\begin{flushleft}
    \Large \textbf{\textcolor{white}{Internship Report : \\ Equidistribution of points on the hypersphere}} \\
    \vspace{0.5cm}
    \normalsize \textcolor{white}{Author : Ly Hai HOANG - 4th-year student in Engineering Physics \\ Internship at Systems Analysis and Architecture Laboratory} \\
    \vspace{0.5cm}
    \normalsize \textcolor{white}{ Supervisors : 
    \begin{itemize}[label=-]
        \item Mr. Iann Geber (LPCNO, INSA Toulouse)
        \item Mr. Antoine Jay (LAAS, SUPAERO)
    \end{itemize}
    \normalsize Internship period : From 02/06/2025 to 29/08/2025}
    
\end{flushleft}

\clearpage

\pagenumbering{arabic} 
\setcounter{page}{1}   

\begin{flushleft}
    \begin{center}
            \Large \textbf{National Institute of Applied Sciences of Toulouse  }
            \includegraphics[width=0.4\linewidth]{image/Logo.png}
    \end{center}
    \Large \textbf {\textcolor{red}{Internship Report : \\ Equidistribution of points on the hypersphere}} \\
    \vspace{2cm}
    \normalsize Author : Ly Hai HOANG \\ 4th-year student of Engineering Physics Department, INSA Toulouse \\ Internship at Systems Analysis and Architecture Laboratory \\
    \vspace{0.5cm}
    \normalsize Supervisors : 
    \begin{itemize}[label=-]
        \item Mr. Iann Geber (LPCNO, INSA Toulouse)
        \item Mr. Antoine Jay (LAAS, SUPAERO)
    \end{itemize}
    \normalsize Internship period : From 02/06/2025 to 29/08/2025
    
\end{flushleft}

\newpage
% ============ concept ============
\tableofcontents
\chapter*{Résumé}
\addcontentsline{toc}{chapter}{\protect\numberline{}Résumé}
\chapter*{Acknowledgements}
\addcontentsline{toc}{chapter}
{\protect\numberline{}Acknowledgements} 
I would like to express my deep gratitude to the Engineering Physics Department at the National Institute of Applied Sciences of Toulouse (INSA Toulouse) for offering me an internship opportunity to conduct research and gain further experience. I would also like to sincerely thank Dr.Andrea Balocchi for allowing me to pursue the double degree between INSA Toulouse and Paul Sabatier University, as well as all the faculty members in the department who accompanied me with the essential knowledge and important skills for a smooth start to my internship.

I would like to thank Dr.Anne Hemeryck, head of the Micro-Nano-Bio Technology (MNBT) Department  at the Scientific Institute for Multi-level Modeling of Materials (m3) Team, for her support with administrative procedures and equipment requirements during my internship at Systems Analysis and Architecture Laboratory (LAAS). I would particularly like to express my sincere and profound gratitude to Dr.Antoine Jay, who supervised me through this internship. Thanks to his enthusiasm, I was able to carry it out as planned and acquire many of the skills that are indispensable in scientific research, as well as all the PhDs, Postdocs, and research interns who supported me during this period. 

\begin{flushright}
Toulouse, \today

Ly Hai HOANG
\end{flushright}


\chapter{Introduction}
The algorithm of equidistribution on a hypersphere created by initialization and optimization methods, the results returned are based on the number of dimensions D and the number of initial points N of the system. This algorithm includes symmetry properties, calculating matrices of vectors and coded by using the Fortran programming language. This topic is part of a larger project aimed at studying the properties of materials at the Multi-level Modeling of Materials Team. This project involves several steps of modeling and simulating phenomena in semiconductor material structures to find compatible defect sites depending on the intended use.This project includes computational simulation of inter-lattice energy, quantum simulation and molecular dynamics simulation. All these steps are taken to predict the possible phenomena in micro nano technology to improve the efficiency of the fabrication and manufacturing process.

This Internship is one of the first steps towards exploring the equidistributed structure of materials based on the uniform distribution on a hypersphere. The equidistribution linked to the study of the potential energy surface (PES) and geometric properties of material's cristalline structure. Energy minima correspond to stable atomic structures, while sadle points represent transition states. The aim of this internship is to develop a routine within the existing exploration code, responsible for generating exploration directions optimally distributed in the hyperspace of atomic positions. The main difficulty lies in taking into account the symmetries of the crystal space group. The method will then be applied to the study of metastable defects in silicon. Finally, a collaborative scientific paper will be written on the theme of equidistribution on a hypersphere.
\chapter{Laboratory presentation}
\section{Systems Analysis and Architecture Laboratory}
\section{Micro-Nano-Bio Technology Department}
\section{Multi-level Modeling of Materials Team}
\chapter{Context general of the Internship}
\chapter{Theorical background}

\section{Definition of the hypersphere}
A hypersphere of radius \(r\) in a Euclidean space of dimension \(D\), denoted  \(\mathbb{S}^{D-1}(r)\) is the set of points located at distance \(r\) from the origin. Formally, the hypersphere is defined by : 
\begin{equation}
    \mathbb{S}^D(r) = \left\{ \mathbf{X} \in \mathbb{R}^{D+1} \ \big| \ \|\mathbf{X}\| = r \right\} = \left\{ \mathbf{X} \in \mathbb{R}^{D+1} \ \big| \ \sum_{i=1}^{D+1} x_i^2  = r^2 \right\} 
\end{equation}
In special cases, when the radius is equal to 1, we call it a unit hypersphere which defined by :
\begin{equation}
    \mathbb{S}^D = \left\{ \mathbf{X} \in \mathbb{R}^{D+1} \ \big| \ \|\mathbf{X}\| = 1 \right\} = \left\{ \mathbf{X} \in \mathbb{R}^{D+1} \ \big| \ \sum_{i=1}^{D+1} x_i^2  = 1 \right\} 
\end{equation}
\begin{figure}[H] 
  \centering
  \includegraphics[width=0.3\linewidth]{image/3-sphere.png}
  \captionsetup{justification=centering}
  \caption{Direct projection of 3-sphere into 3D space and covered with surface grid, showing structure as stack of 3D spheres (2-spheres).}
  %\label{fig:ten_anh}
\end{figure}
The circle is considered 1 dimensional sphere in the 2D plan, and the 2 dimensional sphere is a sphere in the 3D space. In mathematics, there may exist a \(D\) dimensional hypersphere with a non-negative integer D.

\section{Equidistribution of N points on a hypersphere}
The equidistribution of \(N\) points on a hypersphere is simply understood as the smallest distance between nearest neighbor being completely equal and as long as possible, this maximal distance is called the covering radius. The geodesic distance between two vectors \(\mathbf{X_i},\mathbf{X_j} \in \mathbb{S^D}\) on the unit \(D\)-sphere is the angle between these two vectors, is defined by : 
\begin{equation}
    \mu(\mathbf{X_i}, \mathbf{X_j}) = \arccos\left( \mathbf{X_i} \cdot \mathbf{X_j} \right)
\end{equation}
Therefore, the purpose of this subject is to find the smallest angles between adjacent vectors of nearest neighboring vectors.
This subject will focus on finding methods to create and adjust the angles between these vectors to obtain the equidistribution on a hypersphere. By finding the smallest angles between the nearest neigboring vectors, optimization methods will be applied to modify large deviations closer to the average value. Theoretically, the minimum, maximum and average angles between the nearest neighbors are respectively defined by : 
\begin{equation}
\theta_{\min}(A^N) = \min_{\mathbf{X_j} \in A^N} \left[ \min_{\substack{\mathbf{X_i} \in A^N \\ i \ne j}} \mu(\mathbf{X_i},\mathbf{X_j}) \right]
\end{equation}
\begin{equation}
\theta_{\max}(A^N) = \max_{\mathbf{X_j} \in A^N} \left[ \min_{\substack{\mathbf{X_i} \in A^N \\ i \ne j}} \mu(\mathbf{X_i},\mathbf{X_j}) \right]
\end{equation}
\begin{equation}
\langle \theta(A^N) \rangle= \frac{1}{N} \sum_{j=1}^{N} \left[ \min_{\substack{\mathbf{X_i} \in A^N \\ i \ne j}} \mu(\mathbf{X_i},\mathbf{X_j}) \right]
\end{equation}

\section{Platonic solid}
\subsection{Definition of the Platonic solid}
In 3D Euclidean space, a Platonic solid is a convex and regular polyheron that satisfies three conditions following : 
\begin{itemize}
    \item All faces of Platonic solid are the same regular polygon.
    \item All edges of Platonic solid are of equal length.
    \item The platonic solid has the same number of faces meet at each vertex.
\end{itemize}
Only five types of platonic solids exist in nature, including polyhedra, such as : tetrahedron, cube, octahedron, dodecahedron and icosaheron. The Platonic structure were found in all things in nature, such as the crystalline structure of materials and construction of human civilization throughout history. 
\subsection{Combinatorial properties}
Each Platonic solid can be associated with a pair of integers \((p, q)\) which \(p\) represents the number of edges on each face, and \(q\) represents the number of faces that meet at each vertex. This pair is known as the Schläfli symbol and provides a combinatorial description of the polyhedron. Euler's The characteristic formula indicates the relationship between the number of vertices \(V\), edges \(E\), and faces \(F\) as follows:
\begin{equation}
    V-E+F=2
\end{equation}
Each face has \(p\) edges, but each edge is shared by two faces. The total number of edges is :
\begin{equation}
    E = \frac{pF}{2}
\end{equation}
Each vertex is where \(q\) faces meet, and again each edge connects two vertices. Therefore, the total number of edges is consequently : 
\begin{equation}
    E = \frac{qV}{2}
\end{equation}
Since every polyhedron has at least three edges, it also has at least three faces that meet at each vertex, so we have : \( p \geq 3,\ q \geq 3 \). But \(p\) and \(q\) cannot both be greater than 3. Therefore, the only possibilities are \(p=3\) and \(q \geq 3\), or \(q=3\) and \(p \geq 3\). By using 3 equations (4.7), (4.8) and (4.9), we have : 
\begin{equation}
    \frac{1}{q}+\frac{1}{p}=\frac{1}{2}+\frac{1}{E}
\end{equation}
We can solve this equation by combining the conditions behind : \( 3 \leq q < 6 \) with \(p=3\). Obviously, we know that \(q\) is an integer. So \(q\) can only be 3, 4, or 5. Consequently, there are only five possible Platonic solids with five pairs \((p, q)\) and we can calculate the number of faces \(F\), vertices \(V\) and edges \(E\) as a function of \(q\) and \(p\) : 
\begin{equation}
    V=\frac{4p}{4-(p-2)(q-2)}
\end{equation}
\begin{equation}
    E=\frac{2pq}{4-(p-2)(q-2)}
\end{equation}
\begin{equation}
    F=\frac{4q}{4-(p-2)(q-2)}
\end{equation}

\begin{table}[H]
\centering
\begin{tabular}{|c|c|c|c|c|c|c|}
\hline
\rowcolor{red!20}
\textbf{Solid} & \boldmath$p$ & \boldmath$q$ & \textbf{Face Shape} & \textbf{Faces (F)} & \textbf{Vertices (V)} & \textbf{Edges (E)} \\
\hline
Tetrahedron    & 3 & 3 & Triangle & 4  & 4  & 6  \\
\hline
Cube           & 4 & 3 & Square   & 6  & 8  & 12 \\
\hline
Octahedron     & 3 & 4 & Triangle & 8  & 6  & 12 \\
\hline
Dodecahedron   & 5 & 3 & Pentagon & 12 & 20 & 30 \\
\hline
Icosahedron    & 3 & 5 & Triangle & 20 & 12 & 30 \\
\hline
\end{tabular}
\caption{Properties of the five Platonic solids \((p, q)\) with pairs and polyhedral elements.}
\label{tab:platonic_solids}
\end{table}

\subsection{Geometric examples}
In a Platonic solid, all vertices lie on a sphere centered at the origin (circumscribed sphere) and the central angle between two nearest neighbor vectors is the angle formed at the center by the vectors pointing from the origin to each vertex : 
\begin{equation}
    \theta = \arccos\left( \frac{ \mathbf{u} \cdot \mathbf{v} }{ \| \mathbf{u} \| \, \| \mathbf{v} \| } \right)
\end{equation}

\begin{figure}[H] 
  \centering
  \includegraphics[width=0.8\linewidth,page=1]{graphe/platonic1.pdf}
  \captionsetup{justification=centering}
  \caption{Central angle between adjacent vertices for Platonic solids (in degrees):\\ Tetrahedron: 109.47°; \quad Cube: 70.53°; \quad Octahedron: 90°.}
  %\label{fig:ten_anh}
\end{figure}
\begin{figure}[H] 
  \centering
  \includegraphics[width=0.8\linewidth,page=1]{graphe/platonic2.pdf}
  \captionsetup{justification=centering}
  \caption{Central angle between adjacent vertices for Platonic solids (in degrees):\\Dodecahedron: 116.57°; \quad Icosahedron: 63.43°.} 
  %\label{fig:ten_anh}
\end{figure}

\section{Boxplot}




\chapter{Methodology}
\section{Initialization method}
\section{Global optimization method}
\subsection{Tammes's problem}
\subsection{Thomson's problem}
\subsection{Algorithm}
\section{Local optimization method}
\subsection{Method objective}
After using the global optimization method to adjust the equidistribution of \(N\) points on the surface, we need to use the local optimization method to modify the angle between the vectors so that they are as uniform as possible. 

This method will find the smallest angles between any two vectors in N vectors, then adjust the largest and smallest angles in this value set to a critical value. The critical value is the threshold based on the average value and standard deviation of the smallest angles values. To simplify the calculation, it will be based on the dot product of two vectors to derive the distribution of \(\theta\) on the hypersphere. In order to facilitate comprehension, it is imperative to consider the unit hypersphere, wherein the dot product between two vectors is equivalent to the cosine value between the corresponding vectors. 

\subsection{Algorithm}
Two methods exist for modifying the angles : expand the small angles and narrow the large angles. It is important to note that the angles under consideration are all part of the set \(\left\{ \theta \in A^N \right\}\), which contains the smallest possible angles between the vectors. Because the smallest angles have the biggest dot product values, we only need to find the largest dot product value in \(N^2\) value and apply this local method until the requirements are met. In the context of a set \(N\) of vectors distributed on a sphere, the dot product between any two vectors \(\mathbf{X_i}\) and \(\mathbf{X_j}\) can be represented by an \(N \times N\) matrix as follows:
\begin{equation}
    dp = \begin{bmatrix}
    dp_{11} & dp_{12} & \cdots & dp_{1N} \\
    dp_{21} & dp_{22} & \cdots & dp_{2N} \\
    \vdots & \vdots & \ddots & \vdots \\
    dp_{N1} & dp_{N2} & \cdots & dp_{NN}
    \end{bmatrix} =
    \begin{bmatrix}
    \mathbf{X_1} \cdot \mathbf{X_1} & \mathbf{X_1} \cdot \mathbf{X_2} & \cdots & \mathbf{X_1} \cdot \mathbf{X_N} \\
    \mathbf{X_2} \cdot \mathbf{X_1} & \mathbf{X_2} \cdot \mathbf{X_2} & \cdots & \mathbf{X_2} \cdot \mathbf{X_N} \\
    \vdots & \vdots & \ddots & \vdots \\
    \mathbf{X_N} \cdot \mathbf{X_1} & \mathbf{X_N} \cdot \mathbf{X_2} & \cdots & \mathbf{X_N} \cdot \mathbf{X_N}
    \end{bmatrix}
\end{equation}
The algorithm finds the \(N\) largest values, and this set of values is represented by a matrix of size \(1 \times N\), as follows:
\begin{equation}
    dp_{\mathrm{nn}} = \begin{bmatrix}
    dp_{\mathrm{nn}1} & dp_{\mathrm{nn}2} & \cdots & dp_{\mathrm{nn}N}
    \end{bmatrix}
\end{equation}
where \( dp_{\mathrm{nn}i} = \max(dp_{ij}) \) for a fixed row \(i\), with \( j = 1, 2, \ldots, N \). Here, \( j_{\max} \) denotes the column index at which the maximum occurs, and \( i_{\max} \) denotes the corresponding row index. It is important to note that the values of the dot product of the vector itself are not taken into account, as the vector has a value of 1.0. In the algorithm, the values will be set automatically to -1.0 to avoid incorrect maximum values.

\textbf{Maximize the minimum angles}

The objective of this method is focused on maximizing the smallest angles by decreasing the largest dot product values in the set. The vectors $\mathbf{X_i}$ and $\mathbf{X_j}$ are applied to the following formula and normalized :
\begin{equation}
    \mathbf{X_i}\rightarrow \mathbf{X_i} - \alpha \mathbf{X_j}\rightarrow\mathbf{X_i}/||\mathbf{X_i}|| \\
\end{equation}
\begin{equation}
        \mathbf{X_j}\rightarrow \mathbf{X_j} - \alpha \mathbf{X_i}\rightarrow\mathbf{X_j}/||\mathbf{X_j}||
\end{equation}
\begin{figure}[H] 
  \centering
  \includegraphics[width=0.6\linewidth,page=1]{graphe/min_dot_product.pdf}
  \captionsetup{justification=centering}
  \caption{Illustration of the modification to minimize the dot product between two nearest neighbor vector}
  %\label{fig:ten_anh}
\end{figure}
The value of \(\alpha\) is contingent upon the critical value, which is determined through the calculation of the mean value and standard deviation. Specifically, the threshold is often defined as :
\begin{equation}
    dp^+ = \langle dp_{nn}\rangle + \lambda\sigma_{\theta}
\end{equation}

By using the positive solution of second-order equation, we can define the value of \(\alpha\) depend on dot product values of \(dp_t\) and \(dp_{t+1}\) as :
\begin{equation}
    \alpha_{max} = \frac{1 - dp_t \cdot dp_{t+1}}{dp_t-dp_{t+1}} - \sqrt{\left( \frac{1 - dp_t \cdot dp_{t+1}}{dp_t-dp_{t+1}} \right)^2- 1}
\end{equation}
Subsequent to each modification, the value of \(\lambda\) will undergo a reduction, effectively excluding all values that exceed the critical value threshold. As the number of loop iterations increases, it can be expected that the magnitude of all smallest angles will be equal .

\textbf{Minimize the maximal angles}

The objective of this method is focused on minimizing the largest angles in the set of minimum angles by increasing the smallest dot product values in the set. The vectors $\mathbf{X_i}$ and $\mathbf{X_j}$ are applied to the following formula and normalized :
\begin{equation}
    \mathbf{X_i}\rightarrow \mathbf{X_i} + \alpha \mathbf{X_j}\rightarrow\mathbf{X_i}/||\mathbf{X_i}|| \\
\end{equation}
\begin{equation}
    \mathbf{X_j}\rightarrow \mathbf{X_j} + \alpha \mathbf{X_i}\rightarrow\mathbf{X_j}/||\mathbf{X_j}||
\end{equation}
\begin{figure}[H] 
  \centering
  \includegraphics[width=0.6\linewidth,page=1]{graphe/max_dot_product.pdf}
  \captionsetup{justification=centering}
  \caption{Illustration of the modification to maximize the dot product between two nearest neighbor vector}
  %\label{fig:ten_anh}
\end{figure}
Using the same analytical method of maximizing the minimum angles, we find that the value of threshold and \(\alpha\) obtained after each iteration is as follows :
\begin{equation}
    dp^- = \langle dp_{nn}\rangle - \lambda\sigma_{\theta}
\end{equation}
\begin{equation}
    \alpha_{min} = -\frac{1 - dp_t \cdot dp_{t+1}}{dp_t-dp_{t+1}} - \sqrt{\left( \frac{1 - dp_t \cdot dp_{t+1}}{dp_t-dp_{t+1}} \right)^2- 1}
\end{equation}
The value of \(\lambda\) will be reduced according to the maximizing method below, As the number of loop iterations increases, it can be expected that the magnitude of all smallest angles will be
equal.
\newpage
\textbf{Algoritm of the local optimization method}

\begin{figure}[H] 
  \centering
  \includegraphics[width=0.8\linewidth,page=1]{graphe/algo_dot_product.pdf}
  \captionsetup{justification=centering}
  \caption{Algorithm diagram of the dot product modifying method}
  %\label{fig:ten_anh}
\end{figure}




\chapter{Results and Discussion}
\section{Result of initialization method}
\section{Result of global optimization method}
\section{Result of local optimization method}
\section{Comparison between our code and spherical code data}
\chapter{Conclusion and Perspectives}
\chapter{References}


\appendix
\chapter{Conditions of value in the algorithm}
\section{Determining the Coefficient of Displacement Using the Dot Product Modification Method}

In this section, the value of \(\alpha\) is determined by solving a second-order equation that depends on the values of \(dp_t\) and \(dp_{t+1}\). In this algorithm, the value of \(dp_{t+1}\) is selected based on whether it falls within the threshold range \(\langle dp_{nn} \rangle \pm \lambda\sigma_{\theta}\).

For the maximization dot product method, we have : 
\begin{align*}
    & dp_t = \frac{\mathbf{X_i} \cdot \mathbf{X_j}}{||\mathbf{X_i}||.||\mathbf{X_j}||} = \mathbf{X_i} \cdot \mathbf{X_j}\\
    & dp_{t+1} =  \frac{\mathbf{X_i'} \cdot \mathbf{X_j'}}{||\mathbf{X_i'}||.||\mathbf{X_j'}||} = \frac{(\mathbf{X_i} - \alpha \mathbf{X_j})\cdot(\mathbf{X_j} - \alpha \mathbf{X_i})}{(\mathbf{X_i} - \alpha \mathbf{X_j})^2}
\end{align*}
   
    

Where : 
\begin{align*}
    &(\mathbf{X_i} - \alpha \mathbf{X_j}) \cdot (\mathbf{X_j} - \alpha \mathbf{X_i}) \\
    &= \mathbf{X_i} \cdot \mathbf{X_j}
    - \alpha \mathbf{X_i} \cdot \mathbf{X_i}
    - \alpha \mathbf{X_j} \cdot \mathbf{X_j}
    + \alpha^2 \mathbf{X_j} \cdot \mathbf{X_i} \\
    &= \mathbf{X_i} \cdot \mathbf{X_j} - \alpha - \alpha + \alpha^2 \mathbf{X_i} \cdot \mathbf{X_j} \\
    &= (1 + \alpha^2)\mathbf{X_i} \cdot \mathbf{X_j} - 2\alpha \\
    &= (1 + \alpha^2)dp_t - 2\alpha
\end{align*}
And : 
\begin{align*}
    &(\mathbf{X_i} - \alpha \mathbf{X_j}) \cdot (\mathbf{X_i} - \alpha \mathbf{X_j}) \\
    &= \mathbf{X_i} \cdot \mathbf{X_i}
    - \alpha \mathbf{X_i} \cdot \mathbf{X_j}
    - \alpha \mathbf{X_j} \cdot \mathbf{X_i}
    + \alpha^2 \mathbf{X_j} \cdot \mathbf{X_j} \\
    &= 1 - 2\alpha \, dp_t + \alpha^2
\end{align*}
So we have the second-order equation : 
\begin{align*}
    & dp_{t+1} = \frac{(1 + \alpha^2)dp_t - 2\alpha}{1 - 2\alpha \, dp_t + \alpha^2} \\
    &\Rightarrow\; dp_{t+1}(1 - 2\alpha \, dp_t + \alpha^2) = (1 + \alpha^2)dp_t - 2\alpha \\
    &\Rightarrow\; dp_{t+1} - 2\alpha \, dp_{t+1} \, dp_t + \alpha^2 dp_{t+1}
    = dp_t + \alpha^2 dp_t - 2\alpha \\
    &\Rightarrow\; \alpha^2(dp_{t+1} - dp_t) + \alpha(2dp_{t+1}dp_t - 2)
    + (dp_{t+1} - dp_t) = 0 \\
    &\Rightarrow\; \alpha^2 + 2\alpha\left(\frac{dp_{t+1}dp_t - 1}{dp_{t+1} - dp_t} \right) + 1 = 0
\end{align*}

We can solve this equation and choose \(\alpha\) is always a positive solution for this equation : 
\begin{align*}
     \alpha = \frac{1 - dp_t \cdot dp_{t+1}}{dp_t-dp_{t+1}} - \sqrt{\left( \frac{1 - dp_t \cdot dp_{t+1}}{dp_t-dp_{t+1}} \right)^2- 1}
\end{align*}
By solving the minimization of dot product method similarly, we find alpha to be :
\begin{align*}
     \alpha = -\frac{1 - dp_t \cdot dp_{t+1}}{dp_t-dp_{t+1}} - \sqrt{\left( \frac{1 - dp_t \cdot dp_{t+1}}{dp_t-dp_{t+1}} \right)^2- 1}
\end{align*}
\end{document}



